\documentclass[a4paper,11pt]{article}

\usepackage[margin=1.8cm]{geometry}
\usepackage{amsmath,amssymb}
\usepackage{tcolorbox}
\usepackage{xcolor}
\usepackage{enumitem}
\usepackage{titlesec}

% Farben
\definecolor{bsyblue}{HTML}{0066CC}
\definecolor{bsygreen}{HTML}{009966}
\definecolor{bsyred}{HTML}{CC0033}
\definecolor{bsygray}{HTML}{444444}
\definecolor{bsyyellow}{HTML}{FFCC00}
\definecolor{bsybg}{HTML}{F7F9FB}

% Überschriften schöner
\titleformat{\section}{\Large\bfseries\color{bsyblue}}{}{0em}{}
\titleformat{\subsection}{\large\bfseries\color{bsygreen}}{}{0em}{}

\setlist[itemize]{leftmargin=1.2em}
\setlist[enumerate]{leftmargin=1.4em}

\begin{document}

\begin{center}
{\Huge \textbf{BSY Cheatsheet – Julia Edition}}\\[4mm]
{\large \color{bsygray}Alle Aufgabentypen, Formeln und Schritte – idiotensicher.}
\end{center}

%%%%%%%%%%%%%%%%%%%%%%%%%%%%%%%%%%%%%%%%%%%%%%%%%%%%%%%%%%%%
\section{Aufgabentyp 1: Cache \& Memory Access}

\begin{tcolorbox}[colback=bsybg,colframe=bsyblue,title=\textbf{Hitrate bei sequentiellem Zugriff}]
\textbf{Gegeben:} \texttt{int} = 4B, Cacheline = 64B

\[
n_{\text{pro Line}} = \frac{64}{4} = 16
\]

\[
h_c = \frac{16-1}{16} = \frac{15}{16} = 0.9375
\]

\textbf{Merke:} Bei sequentiellem Zugriff ist fast alles ein Hit.
\end{tcolorbox}

\begin{tcolorbox}[colback=bsybg,colframe=bsygreen,title=\textbf{AMAT – mittlere Speicherzugriffszeit}]

\[
T_a = h_1 T_1 + (1-h_1)\bigl(h_2 T_2 + (1-h_2)(h_3 T_3 + (1-h_3)T_{\text{mem}})\bigr)
\]

\textbf{Rezept:}
\begin{enumerate}
\item Zyklen in ns umrechnen.
\item Formel von innen nach außen ausrechnen.
\item Ergebnis in ns angeben.
\end{enumerate}
\end{tcolorbox}

%%%%%%%%%%%%%%%%%%%%%%%%%%%%%%%%%%%%%%%%%%%%%%%%%%%%%%%%%%%%
\section{Aufgabentyp 2: Makefile}

\begin{tcolorbox}[colback=bsybg,colframe=bsyblue,title=\textbf{Makefile Regeln – Minimalwissen}]
\begin{itemize}
\item \texttt{\$@} = Target
\item \texttt{\$<} = erste Dependency
\item \texttt{\$^} = alle Dependencies
\item \textbf{.c → .o} wenn .c neuer ist als .o
\item Executable steht in der Link-Zeile: \texttt{-o \$@.e}
\end{itemize}
\end{tcolorbox}

%%%%%%%%%%%%%%%%%%%%%%%%%%%%%%%%%%%%%%%%%%%%%%%%%%%%%%%%%%%%
\section{Aufgabentyp 3: \texttt{fork()} – Prozessbäume}

\begin{tcolorbox}[colback=bsybg,colframe=bsyred,title=\textbf{Rezept: Anzahl Prozesse}]
\begin{enumerate}
\item Prozessbaum zeichnen.
\item Jeder \texttt{fork()} verdoppelt.
\item Anzahl Blätter = Anzahl Prozesse.
\end{enumerate}
\end{tcolorbox}

\begin{tcolorbox}[colback=bsybg,colframe=bsygreen,title=\textbf{Rezept: Anzahl Ausgaben}]
\begin{enumerate}
\item Baum zeichnen.
\item Prüfen, welche Pfade die Ausgabe erreichen.
\item Pro Endprozess zählen.
\end{enumerate}
\end{tcolorbox}

%%%%%%%%%%%%%%%%%%%%%%%%%%%%%%%%%%%%%%%%%%%%%%%%%%%%%%%%%%%%
\section{Aufgabentyp 4: Buddy System}

\begin{tcolorbox}[colback=bsybg,colframe=bsyblue,title=\textbf{Rezept: interne Fragmentierung}]
\begin{enumerate}
\item Für jede Anfrage \(S\): kleinstes \(2^k\) suchen.
\item Verlust: \(V = 2^k - S\)
\item Gesamtverlust: Summe aller \(V\)
\end{enumerate}
\end{tcolorbox}

%%%%%%%%%%%%%%%%%%%%%%%%%%%%%%%%%%%%%%%%%%%%%%%%%%%%%%%%%%%%
\section{Aufgabentyp 5: Paging}

\begin{tcolorbox}[colback=bsybg,colframe=bsygreen,title=\textbf{Wichtigste Formeln}]

\[
\text{Page Size} = 2^{\text{OffsetBits}}
\]

\[
N_{\text{PD}} = 2^{\text{DirectoryBits}}
\]

\[
N_{\text{PT}} = 2^{\text{PageNumberBits}}
\]

\[
\text{Size}_{\text{PD}} = N_{\text{PD}} \cdot \text{PointerSize}
\]

\end{tcolorbox}

%%%%%%%%%%%%%%%%%%%%%%%%%%%%%%%%%%%%%%%%%%%%%%%%%%%%%%%%%%%%
\section{Aufgabentyp 6: Page Replacement (LRU)}

\begin{tcolorbox}[colback=bsybg,colframe=bsyred,title=\textbf{LRU – idiotensicheres Rezept}]
\begin{enumerate}
\item Frames leer initialisieren.
\item Für jede Page:
\begin{itemize}
\item Hit → Reihenfolge aktualisieren.
\item Miss → Page Fault markieren.
\item Freier Frame? → einfügen.
\item Sonst → \textbf{LRU ersetzen}.
\end{itemize}
\item Bei Gleichstand → Frame mit kleinster Nummer.
\end{enumerate}
\end{tcolorbox}

%%%%%%%%%%%%%%%%%%%%%%%%%%%%%%%%%%%%%%%%%%%%%%%%%%%%%%%%%%%%
\section{Aufgabentyp 7: Scheduling}

\subsection*{Real-Time (EDF)}

\begin{tcolorbox}[colback=bsybg,colframe=bsyblue,title=\textbf{EDF Rezept}]
\begin{enumerate}
\item Deadlines berechnen: \(D_{i,k} = kT_i\)
\item Zu jedem Rescheduling:
\begin{itemize}
\item Task mit frühester Deadline wählen.
\item Bei Gleichstand: kürzere Periode gewinnt.
\end{itemize}
\item Gantt-Diagramm zeichnen.
\end{enumerate}
\end{tcolorbox}

\subsection*{Multilevel + RR(q=1)}

\begin{tcolorbox}[colback=bsybg,colframe=bsygreen,title=\textbf{RR + Prioritäten Rezept}]
\begin{enumerate}
\item Zeitachse in 1er-Schritten.
\item Neue Prozesse einfügen.
\item Höchste Priorität wählen.
\item Innerhalb: Round Robin mit \(q=1\).
\item Burst--1, bei 0 → Prozess fertig.
\end{enumerate}
\end{tcolorbox}

%%%%%%%%%%%%%%%%%%%%%%%%%%%%%%%%%%%%%%%%%%%%%%%%%%%%%%%%%%%%
\section{Aufgabentyp 8: Theorie}

\begin{tcolorbox}[colback=bsybg,colframe=bsyyellow,title=\textbf{Mini-Merkzettel}]
\begin{itemize}
\item System Calls = Kernel-Funktionen (I/O, Sleep, Prozesse)
\item PCB = OS-Datenstruktur (Status, Register, PID)
\item Zombie = Prozess beendet, Elternprozess hat nicht gewartet
\item Paging = Pages und Frames \textbf{gleich gross}
\item Interne Fragmentierung = Block zu groß
\item Externe Fragmentierung = viele kleine Lücken
\end{itemize}
\end{tcolorbox}

\end{document}
